\chapter[Time responses]{Time responses without rational approximation}\label{ch:response}

The previous chapters located poles and zeros. This chapter computes
\emph{time responses} --- impulse, step and forced --- of the same
nonrational transfer functions, with the functions \code{cimpulse},
\code{cstep}, \code{clsim} and \code{cinvlaplace}.

\section{Motivation}

For a rational model, MATLAB's \code{impulse}, \code{step} and \code{lsim}
integrate a state-space realization in time. A transfer function with a
delay, a diffusion term or unsteady aerodynamics has no finite realization,
and the usual practice simulates a rational approximation instead: a Padé
model of the delay, a truncated modal expansion of the PDE, or aerodynamic
lag states. Chapter~\ref{ch:delay} showed that such approximations can
change stability conclusions; they can change transients as well.

The response of a causal linear system at rest is the inverse Laplace
transform of $G(s)U(s)$,
\begin{equation}
  y(t)=\frac{1}{2\pi\ii}\int_{\sigma-\ii\infty}^{\sigma+\ii\infty}
        e^{st}\,G(s)U(s)\,ds ,
  \label{eq:bromwich}
\end{equation}
where the Bromwich line $\Real s=\sigma$ lies to the right of every
singularity of $G(s)U(s)$. Equation~\eqref{eq:bromwich} involves only values
of $G$: it can be evaluated numerically for any transfer function that can
be evaluated. What is discretized is this integral, not the physical model,
and the discretization is refined until it converges.

\section{Problem statement and the domain assumption}

The module computes zero-state responses of real, causal, single-input
single-output systems: all internal states, including the memory of a
delay line, of a PDE or of an aerodynamic wake, are zero at $t=0$. A
transfer function does not describe an arbitrary internal initial state,
so no analogue of MATLAB's \code{initial} is provided.

The one piece of information that cannot be extracted from a function
handle is where $G$ is analytic. The user supplies it as
\code{SingularityBound}~$a$: the assertion that $G$ is analytic for
$\Real s>a$ (or directly as the line, \code{Abscissa}). For coefficient
vectors and \code{tf}/\code{zpk}/\code{ss} models it is computed from the
poles. Table~\ref{tab:bounds} gives typical justifications.

\begin{table}[H]\centering\small
\caption{Justifying the inversion half-plane.}\label{tab:bounds}
\begin{tabularx}{\textwidth}{>{\raggedright}p{4.2cm}p{1.1cm}X}
\toprule
Transfer function & Bound & Justification\\\midrule
$1/(s+1+0.5e^{-s})$ & 0 & for $\Real s\ge0$, $|s+1|\ge1>0.5\ge|0.5e^{-s}|$: no zeros of the denominator\\
$e^{-a\sqrt s}$ (diffusion) & 0 & only singularity: the branch cut $(-\infty,0]$\\
$e^{-\tau s}/(s+1)$ & $-1$ & pole at $-1$; $e^{-\tau s}$ is entire\\
unstable system & growth rate & the unstable poles must lie left of the line\\
aeroelastic section, 180 ft/s & 5 & the right-half-plane count of Chapter~\ref{ch:aero} gives the pair $2.13\pm74.8\ii$\\
\bottomrule\end{tabularx}
\end{table}

A bounded root search helps to choose the bound but cannot prove it: a
singularity outside the searched rectangle would be missed. The bound is
therefore recorded as an assertion (\code{info.domainSource}). Choosing a
line to the left of an unstable pole does not fail loudly; it computes a
different, non-causal function (Figure~\ref{fig:respconv}, right).

\section{Inversion on a Bromwich line}

\paragraph{Discretization.} On $s=\sigma+\ii\omega$, \eqref{eq:bromwich}
is the Fourier integral of $f_\sigma(t)=e^{-\sigma t}f(t)$, with $F=GU$.
Sampling $F$ at $\omega_k=2\pi k/P$ and applying the inverse FFT of length
$N$ gives, at $t_n=nP/N$,
\begin{equation}
  \frac{1}{\Delta t}\operatorname{IFFT}\bigl(F(\sigma+\ii\omega_k)\bigr)_n
  =\sum_{m\ge0} f_\sigma(t_n+mP) + E_{\mathrm{band}},
  \qquad \Delta t=\frac{P}{N},
  \label{eq:aliasing}
\end{equation}
where the sum is the periodic repetition of $f_\sigma$ (aliasing in time;
only $m\ge0$ appears because $f$ is causal) and $E_{\mathrm{band}}$ is the
error of truncating the frequency band at $\pm\pi N/P$. Multiplying by
$e^{\sigma t_n}$ recovers $f(t_n)$ plus the aliasing terms
$e^{-m\sigma P}f(t_n+mP)$, $m\ge1$. If $|f(t)|\lesssim Ce^{at}$, the first
of them is of relative order $e^{-(\sigma-a)P}$. The solver chooses
\[
  \sigma=\max(a,0)+\frac{\max\{12,\;-\ln(0.01\,\min(\mathrm{AbsTol},\mathrm{RelTol}))\}}{P},
\]
which makes the aliasing contribution smaller than the requested tolerance
by a factor of about 100. The price is the amplification $e^{\sigma t}$ of
rounding errors when the weighting is undone, reported in
\code{info.amplification}; very long horizons for growing responses are
refused.

\paragraph{Refinement.} Each sample is computed four times: with the
reference period and resolution, with the period doubled, with the
bandwidth doubled, and with the line shifted to the right by $0.5/P$. The
three differences --- period, bandwidth and shift --- estimate the three
error sources separately. A sample is accepted when they are below
$\mathrm{AbsTol}+\mathrm{RelTol}\,|y|$ in two successive rounds;
otherwise period and bandwidth are refined. Figure~\ref{fig:respconv}
(left) shows the three differences for a delayed first-order system.
Zero-padding a previously sampled spectrum would add no information; each
refinement evaluates $G$ at new frequencies.

\begin{figure}[H]\centering
\includegraphics[width=\textwidth]{time_response_convergence.png}
\caption{Left: period, bandwidth and shift differences during refinement
($e^{-0.713s}/(s+1)$). Right: for the unstable $1/(s-0.3)$, inversion on a
line right of the pole gives the causal response $e^{0.3t}$; an inverse FFT
on the imaginary axis gives a different function.}
\label{fig:respconv}
\end{figure}

\paragraph{Independent methods.} Two further methods accept arbitrary
positive times and serve as cross-checks. \code{Method='dehoog'} sums the
Fourier series of \eqref{eq:aliasing} with the continued-fraction
acceleration of \citet{dehoog1982}; its quotient-difference table is a
Padé-type device for summing the series, not a rational approximation of
$G$. It checks the degree, the period and the line shift; each degree
check doubles the degree, hence the sampled bandwidth, so that consecutive
degrees cannot share a plateau that omits lightly damped poles just above
the band (Chapter~\ref{ch:wing}).
\code{Method='quadrature'} integrates a truncated Bromwich integral with
adaptive Gauss--Kronrod quadrature and checks the cutoff, the quadrature
tolerance and the line; it avoids the conditioning limits of the
quotient-difference table in double precision for late oscillatory
responses. There is no automatic switching between methods.

\section{Dirac terms, the origin and jumps}

\paragraph{Direct terms.} If $G=D+G_r$ with a constant $D$, the impulse
response contains $D\delta(t)$ and, with a transport delay $\tau$ of an LTI
model, $D\delta(t-\tau)$. A vector of samples cannot represent these; the
returned impulse response is the ordinary part $g_r$, and the Dirac terms
are listed in \code{info.singularTerms}. For function handles no hidden
Dirac term can be excluded from finite samples: \code{cimpulse} requires
the assertion \code{RegularImpulse=true}, with an optional known
\code{Feedthrough}.

\paragraph{The origin.} An impulse response with $g(0^+)=h_0\neq0$ jumps at
$t=0$, which would slow the convergence of the inversion. When $h_0$ is
known --- exactly for rational models, or through \code{InitialValue} ---
the solver inverts $G_r(s)-h_0/(s+\beta)$, whose inverse is continuous at
zero, and adds $h_0e^{-\beta t}$ back exactly; $\beta>0$ is chosen so that
the auxiliary pole lies left of the line. When $h_0$ is unknown, the value
at $t=0$ is returned as \code{NaN} and marked unresolved; positive times are
unaffected. Kernels that are infinite at zero, such as $1/\sqrt{\pi t}$ for
$G=s^{-1/2}$, are computed at positive times.

\paragraph{Jumps.} At a jump of $f$ the inversion integral converges to the
midpoint of the jump. Known jumps --- the direct term of a step, the delay
of an LTI model --- are placed exactly with the right-continuous
convention. A jump unknown to the solver, such as a delay inside a function
handle, leaves the single sample at that time unresolved, while the other
samples converge.

\section{Forced response and input reconstruction}

\code{clsim} computes the response to an input given by samples $u_j$ at
uniform times $t_j$ (or a handle sampled there). The sampled input must be
completed between samples; the user chooses a zero-order hold (ZOH) or a
first-order hold (FOH), and the response is computed exactly for that
reconstructed input. Convolving sampled impulse responses would require
$g(0)$, which may be infinite; instead the solver uses integrated kernels.
With $G=D+G_r$, define the step and ramp responses
$S_r=\mathcal L^{-1}\{G_r/s\}$ and $R_r=\mathcal L^{-1}\{G_r/s^2\}$, both
zero before $t=0$. A held input is a sum of delayed steps, one per jump
$u_j-u_{j-1}$ (with $u_{-1}=0$):
\begin{equation}\label{eq:zoh-sum}
  y(t)=Du(t)+\sum_{t_j\le t}(u_j-u_{j-1})\,S_r(t-t_j).
\end{equation}
A piecewise-linear input with slopes $m_j=(u_{j+1}-u_j)/(t_{j+1}-t_j)$ is a
step $u_0$ plus delayed ramps, one per change of slope:
\begin{equation}
  y(t)=Du(t)+u_0S_r(t)+m_0R_r(t)+\sum_{j\ge1,\,t_j\le t}(m_j-m_{j-1})\,R_r(t-t_j).
\end{equation}
On a uniform grid these sums are discrete linear convolutions. They are
computed with zero-padded FFTs (length at least $2n-1$, so that the
convolution is linear, not circular), after the exact rescaling
$\sum_j w_jK(t_n-t_j)=e^{\sigma t_n}\sum_j (w_je^{-\sigma t_j})
(K(t_n-t_j)e^{-\sigma(t_n-t_j)})$, which limits the dynamic range for
growing kernels. The kernel error estimates are propagated through the same
sums. For rational models the result agrees with MATLAB's \code{lsim} with
the same hold to about $10^{-10}$ (ZOH) and $10^{-9}$ (FOH). With an
\code{InputDelay} $T$ and $u_0\ne0$, \code{lsim}'s FOH lets the delayed input
rise linearly from zero over the sample interval ending at $T$, while
\code{clsim} keeps the jump of the interpolant at $T$; the two then differ
by about $g(0^+)\,u_0\,\Delta t/2$ just after $T$. The delay-free response
shifted by $T$ confirms \code{clsim} to $10^{-8}$.

\subsection{Several inputs with the same model}

In \eqref{eq:zoh-sum} and its FOH analogue, only the weights depend on the
input; the kernels $S$ and $R$ depend on the system alone. When many inputs
are applied to the same system, such as a sine sweep, a family of reference
signals or random disturbances, the inversions can be done once.
\code{ckernel} performs them and returns a read-only object \code{K};
\code{clsim} accepts \code{K} in place of $G$ and only evaluates the
convolution. The example below ``measures'' the frequency response of the
delayed loop by simulation, fitting a sinusoid to the steady part of each
output:
\begin{lstlisting}
G = @(s) 1./(s + 1 + 0.5*exp(-s));
t = (0:0.02:20).';
K = ckernel(G, t, 'SingularityBound', 0);    % inversions, once
w = 0.5:0.25:2;  late = t > 10;
for k = 1:numel(w)
    [y,~,info] = clsim(K, sin(w(k)*t), t);   % convolution only
    ab = [sin(w(k)*t(late)) cos(w(k)*t(late))] \ y(late);
    gain(k) = norm(ab);                      % compare with abs(G(1i*w(k)))
end
\end{lstlisting}
The seven gains agree with $|G(i\omega)|$ to about $10^{-4}$, the size of
the transient left at $t = 10$ (the rightmost characteristic root has
$\operatorname{Re}s\approx-1.10$). Every \code{info.evaluations} is zero,
and the script takes about 1.2\,s against 2.8\,s for seven ordinary calls.

\paragraph{What is stored.} FOH preparation stores $R_r$ and $S_r$, so a
later input may start at $u_0\ne0$; ZOH stores $S_r$ only. The direct term
and a known transport delay are kept. \code{K} holds numbers for one model,
one grid and one hold, not the function handle: changing a parameter, the
time samples or the hold requires a new preparation, and \code{clsim}
rejects a mismatched grid, hold or model option. Each call is an
independent zero-state experiment.

\paragraph{Two tolerances.} The error of the stored kernels is amplified by
each input. With a step-kernel error estimate $\epsilon_S$, the ZOH output
error is estimated as
\begin{equation}
 \epsilon_y(t_n)=\sum_{j=0}^{n}|u_j-u_{j-1}|\,\epsilon_S(t_n-t_j),
\end{equation}
and FOH propagates $\epsilon_R$ through the slope changes and adds
$|u_0|\,\epsilon_S$. The kernels are therefore prepared at tighter
tolerances ($10^{-8}$ absolute, $10^{-6}$ relative) than the output
($10^{-6}$, $10^{-4}$), and $\epsilon_y$ is recomputed and checked for every
input. A large or rapidly varying input may fail that check even though the
preparation converged; the samples are then marked unresolved and the
kernels are \emph{not} refined automatically. The user either prepares
tighter kernels or relaxes an unnecessarily strict output tolerance. The
estimate is conservative (true errors near $10^{-9}$ when it first exceeds
$10^{-6}$ in the tests), and it retains the finite-bandwidth caveat of the
inversion: it is not a certified bound.

\section{Validation}

The study \path{examples/time_response/run_time_response_study.m} compares
step responses computed from transfer functions with references obtained by
different means (Table~\ref{tab:respval}, Figure~\ref{fig:respval}). The
finite-element and modal truncations belong to the references, not to the
inversion.

\begin{table}[H]\centering\small
\caption{Maximum absolute errors of step responses on $0\le t\le6$.}\label{tab:respval}
\begin{tabular}{llr}\toprule
Transfer function & Independent reference & Error\\\midrule
$1/(s+1+0.5e^{-s})$ & method-of-steps series & $1.2\times10^{-8}$\\
$e^{-0.7\sqrt s}$ & $\operatorname{erfc}(0.7/(2\sqrt t))$ & $4.6\times10^{-8}$\\
heated rod (Chapter~\ref{ch:distributed}) & separation of variables, 150 modes & $1.1\times10^{-7}$\\
beam and oscillator (Chapter~\ref{ch:beam}) & finite elements, 32 elements & $2.5\times10^{-5}$\\
aeroelastic pitch, 160 and 180 ft/s & quadrature versus FFT & $<10^{-10}$\\
\bottomrule\end{tabular}
\end{table}

\begin{figure}[H]\centering
\includegraphics[width=\textwidth]{time_response_validation.png}
\caption{Step responses computed from the transfer functions (solid) and
independent references (dashed).}
\label{fig:respval}
\end{figure}

The aeroelastic example of Chapter~\ref{ch:aero} gives the pitch response to
a pitch moment: the transfer function is the second entry of
$D(s,U)^{-1}[0,1]^T$ --- an input/output channel, not the characteristic
determinant --- with exact Theodorsen aerodynamics
(Figure~\ref{fig:respaero}). Below flutter the oscillation decays to the
static deflection; above flutter it grows at the rate predicted by the
right-half-plane count.

\begin{figure}[H]\centering
\includegraphics[width=\textwidth]{time_response_aeroelasticity.png}
\caption{Pitch response to a unit pitch-moment step of the NASA section,
below and above flutter. Circles: independent Bromwich quadrature.}
\label{fig:respaero}
\end{figure}

The regression tests add analytical first-order, delayed, unstable,
marginal and fractional cases, MATLAB's \code{lsim} with both holds, direct
terms, unknown initial limits, scalar-only evaluators and resource limits.

\section{Limitations}

\begin{itemize}
  \item Zero-state, real, single-input single-output responses; for matrix
        models, a chosen input/output channel.
  \item The inversion half-plane is an assertion for function handles.
  \item The input is the chosen ZOH or FOH interpolant of its samples.
  \item Jumps and corners unknown to the solver leave the sample at that
        time unresolved.
  \item Convergence is numerical (successive refinements agree); a feature
        outside the sampled band cannot be excluded, and
        \code{info.certified} is always false.
\end{itemize}
